\documentclass[11pt]{article}
\usepackage[margin=1in]{geometry}
\usepackage{amsmath,amssymb,amsthm}
\usepackage[hidelinks]{hyperref}

\newcommand{\R}{\mathbb{R}}
\newcommand{\T}{\mathsf{T}}
\newcommand{\cH}{\{-1,1\}^n}

\title{Literature-Implied Negative Answer to the Average-Distortion Extension}
\author{arXiv:2310.01868 target; arXiv:2410.21931 support}
\date{2026-07-03}

\begin{document}
\maketitle

\paragraph{Status.}
\emph{Literature-implied answer, negative for the dimension-refined extension in Discussion item 3 of arXiv:2310.01868.}
The relation is an agent-identified implication of the proposition labelled \texttt{prop:Rn more genral} in arXiv:2410.21931; the supporting authors do not appear to explicitly cite or answer this discussion item.

\paragraph{Source question.}
In Section ``Discussion and open problems,'' item 3, Eskenazis asks whether the bound of Theorem labelled \texttt{thm:main} can be extended from bi-Lipschitz embeddings to average-distortion embeddings.  The theorem gives, for finite-dimensional normed \(X\),
\[
  \mathsf{c}_X(\{-1,1\}^n)
  \gtrsim
  \sup_{1\le p\le 2}
  \frac{n}{\T_p(X)\min\{n,\dim(X)\}^{1/p}}.
\]
The same discussion notes that IVV-type arguments imply only the dimension-free average-distortion lower bound
\[
  \T_p(X)^{-1}n^{1-1/p}.
\]

\paragraph{Supporting theorem.}
Chang--Naor--Ren, arXiv:2410.21931, Proposition labelled \texttt{prop:Rn more genral}, states that for \(1\le q\le 2\), the smallest \(D\ge 1\) for which \(\ell_q^n\) embeds into \(\R\) with \(p\)-average distortion \(D\) is, up to universal constants,
\[
  \sqrt{\max\{1,n/p\}}.
\]
Their definition of \(p\)-average distortion is the standard Rabinovich--Naor one: for every Borel probability measure on the source, there is a Lipschitz map satisfying the required \(L_p\)-average lower separation after scaling.

\paragraph{Implication for the source question.}
Take \(X=\R\).  Then \(\dim X=1\) and \(\T_2(\R)=1\).  The \(p=2\) term in the source theorem bound is therefore of order
\[
  \frac{n}{\T_2(\R)\min\{n,1\}^{1/2}} = n.
\]
If that dimension-refined lower bound extended to quadratic average distortion, then \(\cH\) would require quadratic average distortion \(\gtrsim n\) into \(\R\).

But the supporting proposition with \(q=1\) and \(p=2\) gives an embedding of \(\ell_1^n\) into \(\R\) with quadratic average distortion \(O(\sqrt n)\).  Restricting that embedding to the subset \(\cH\) preserves the same average-distortion upper bound for every probability measure on \(\cH\), by extending the measure by zero to \(\ell_1^n\).  On \(\cH\), the inherited \(\ell_1\) metric is exactly twice the Hamming metric, and scaling the source metric changes neither the order nor the validity of an average-distortion estimate.

Consequently,
\[
  \mathrm{Av}_2(\cH,\R) \lesssim \sqrt n,
\]
whereas the proposed dimension-refined extension would force
\[
  \mathrm{Av}_2(\cH,\R) \gtrsim n.
\]
For large \(n\), these are incompatible.  Thus the source theorem's dimension-refined lower bound does not extend to average-distortion embeddings in this form.

\paragraph{Scope and limitations.}
This negative answer targets the low-dimensional improvement in Theorem labelled \texttt{thm:main}.  It does not refute the IVV-style average lower bound quoted in the source discussion: for \(X=\R\) and \(p=2\), that bound is order \(\sqrt n\), which is consistent with the supporting proposition.

\begin{thebibliography}{9}
\bibitem{Eskenazis2310}
Alexandros Eskenazis.
\newblock \emph{Some geometric applications of the discrete heat flow}.
\newblock arXiv:2310.01868, 2023.

\bibitem{CNR2410}
Alan Chang, Assaf Naor, and Kevin Ren.
\newblock \emph{Random zero sets with local growth guarantees}.
\newblock arXiv:2410.21931, 2024.
\end{thebibliography}

\end{document}
